\raggedright
\chapter*{A Note to the Reader}
\addcontentsline{toc}{chapter}{A Note to the Reader}

This is the short book of \emph{Adversarial Cooperation}. Each chapter holds a
thought we can return to: a question about trust, a limit on power, a way of
cooperating without surrendering every secret.

Read these pages slowly, or begin with the question nearest to your own life.
Their brevity leaves room for reflection. It also leaves work to be done.
Some chapters describe mechanisms already studied in the companion; others
ask questions whose answers we do not yet have.

The \emph{Research Companion} follows the same twenty-two chapters and five
appendices. It holds the models, arguments, algorithms, citations, and links
to executable demonstrations where they exist. Each short chapter points to
its corresponding place there. A question that remains open in the companion
remains open here.

Keep both PDFs in the same folder to follow the chapter links. If a reader
does not support links between PDFs, the matching chapter number or appendix
letter provides the same route.

Peace gives this work its purpose. Each mechanism must answer for the smaller
promise it can actually keep.

\vfill
{\small\color{quiet}First reading-edition draft.\par
The short book and its companion are meant to grow together.}
